\section{Institutional Background}
\label{sec:institutional}

The institutional setting we exploit has three layers: a constitutional
right to health that drives mass litigation; a mandatory electronic
procurement platform (BEC) that records every transaction at bid level;
and a parallel administrative-request channel that, since
\BPsampleStartYear, has run alongside the litigated channel under
identical planning constraints but \emph{without} personal-sanction
exposure for the procurement officer. The third layer is unusual---to
our knowledge, no other Brazilian state and no Latin American
jurisdiction we have surveyed operates a comparable dual channel for
urgent pharmaceutical demand. We describe each layer and document the
key institutional fact that identifies our comparison: litigated and
administrative urgent purchases of the same item, on the same platform,
under the same auction procedures, differ only in the penalty regime
faced by the official.

\subsection{Health as a Constitutional Right and the Rise of Litigation}
\label{sec:health_right}

Brazil's 1988 Constitution declares health ``a right of all and a duty
of the state'' (Article~196). The 1988 framework set up the Unified
Health System (SUS), which maintains a formulary of medications subject
to cost-effectiveness criteria \citep{castro2019, soares2019}. Brazilian
courts have adopted a strict literal interpretation of Article~196,
under which any individual may obtain virtually any medication through
litigation---from common analgesics to rare-disease treatments costing
over US\$\BPrareDrugCostUSD{} per year. Two consolidated lines of
Supreme Federal Court (STF) jurisprudence anchor the regime: federal,
state, and municipal governments are jointly liable for medication
delivery once a court order issues, and lower courts retain wide
discretion to grant preliminary injunctions before merits review.\footnote{Key
STF rulings: Recurso Extraordin\'ario 566.471/RN (Tema 6, 2020) and
Recurso Extraordin\'ario 657.718/MG (Tema 500, 2019), establishing the
joint-liability and high-cost-medication doctrines respectively.}

Access to the courts is inexpensive---a prescription and a public
defender suffice---and approximately \BPlitGrantedShare{} of
first-instance claims are granted \citep{cnj2019judicializacao}. The
result is mass health litigation: \BPlitClaimsBR{} first-instance claims
nationally in \BPcnjClaimsYear, with S\~ao Paulo alone recording a
\BPspGrowthSince{} increase between \BPspGrowthBaseYear{} and
\BPsampleEndYear{} (about $22\%$ compounded annually; Online Appendix,
Figure~\ref{fig:distribution}). Litigated health spending in S\~ao Paulo
amounts to approximately US\$\BPlitSpendingUSDmm~million annually---about
\BPlitSpendingShare{} of the state public health budget.

\subsection{The Sanction Channel: Litigated vs.\ Administrative Urgent Purchases}
\label{sec:sanction_channel}

The S\~ao Paulo State Department of Health (SES/SP) delivers health
services through \BPpbuCount{} decentralized public buyer units (PBUs).
All purchases---ordinary and urgent---are conducted through the
\textit{Bolsa Eletr\^onica de Compras} (BEC), the state's electronic
procurement platform, mandatory for common goods since
\BPbecMandatorySinceYear. Two competitive procedures are used on BEC:
\textit{convite}, a sealed-bid tendering procedure used for smaller
purchases, and \textit{preg\~ao eletr\^onico}, a multi-round descending
auction used for larger purchases and the dominant procedure across
urgent transactions in our sample.\footnote{Among the both-types cells
that identify our tightest specifications, \BPbothCellModalityShareBoth{}
of observations are processed through preg\~ao eletr\^onico.} For
ordinary purchases, the internal planning phase typically spans
\BPordPlanLow--\BPordPlanHigh~days, allowing the buyer to aggregate
demand across PBUs and run competitive auctions on relaxed
timelines. \BPinjunctionShare{} of court decisions in our sample arrive
as preliminary injunctions with delivery deadlines of
\BPurgentPlanLow--\BPurgentPlanHigh~days, compressing the internal phase
to roughly one-third of the ordinary timeline.

The penalties for non-compliance with a court order are concrete and
personal. The most common sanction is a daily fine
(\textit{astreinte}), set in the injunction itself; calibrated to the
purchase value and frequently reaching the order value within weeks
of accumulation, the daily fine is the binding constraint procurement
officers report internally. Additional sanctions include administrative
and criminal liability for the responsible official and sequestration
or blocking of state funds. These penalties are public information,
since the court order must be referenced in the tender notice;
suppliers see the sanction context as they bid.

The institutional feature that identifies our comparison is the
\textit{administrative-request channel}, instituted by the SES/SP in
\BPsampleStartYear. The channel allows patients (or, more often,
caregivers and SUS social workers acting on the patient's behalf) to
request a medication directly through an SES/SP intake before any
litigation occurs. Requests are evaluated by a scientific committee
under evidence-based-medicine and cost-effectiveness criteria; admitted
requests proceed through urgent procurement on BEC under the same
compressed-timeline regime as litigated purchases. The relevant
asymmetry is in the penalty regime: \textbf{failure to complete an
administrative purchase carries no court order, no daily fine, no fund
seizure, and no personal liability for the procurement officer.} The
supplier pool is identical: among urgent winners, \BPsupFirmsBothShare{}
also serve ordinary procurement of the same item-class, so the
administrative regime is not characterized by a separate set of
suppliers.

Between \BPsampleStartYear{} and \BPsampleEndYear, the administrative
channel generated approximately \BPadminPurchasesTotal{} purchases,
or about \BPadminToCourtShare{} of the litigated volume. The smaller
absolute count reflects the committee's gating role and the resource
constraint on its capacity, not a difference in the procurement
procedure once a case is admitted. We exploit this institutional
asymmetry as a within-item, within-buyer comparison and discuss the
selection-into-admin wedge in Section~\ref{sec:empirical}, where we
bound it formally with Manski-Lee monotone restrictions and a Heckman
parametric correction.

Figure~\ref{fig:purchase_types} summarizes the three purchase types.

\begin{table}[ht]
  \centering
  \caption{Purchase Types: Institutional Characteristics}
  \label{fig:purchase_types}
  \small
  \begin{threeparttable}
  \begin{tabular}{lccc}
    \toprule
    & Ordinary & Administrative & Litigated \\
    \midrule
    Trigger              & Standard demand & Patient request & Court order \\
    Gatekeeper           & PBU procurement plan & SES/SP scientific committee & Lower-court judge \\
    Source of funds      & Dedicated budget & Diverted budget & Diverted budget \\
    Quantity             & Large            & Small            & Small \\
    Delivery deadline    & \BPordPlanLow--\BPordPlanHigh{} d & \BPurgentPlanLow--\BPurgentPlanHigh{} d & \BPurgentPlanLow--\BPurgentPlanHigh{} d \\
    Penalty for failure  & None             & None             & \textbf{Daily fine, liability, fund seizure} \\
    Visible to bidders   & ---              & ---              & Court order in tender notice \\
    Active since         & Pre-\BPbecMandatorySinceYear & \BPsampleStartYear & 1990s \\
    \bottomrule
  \end{tabular}
  \begin{tablenotes}
    \small
    \item \textit{Notes:} Administrative and litigated purchases share
    all execution constraints (compressed timeline, small quantity,
    diverted budget, BEC platform, identical auction procedures) and
    differ only in the penalty regime visible to the procurement
    officer and to bidders. The within-urgent contrast we exploit in
    Section~\ref{sec:empirical} compares these two columns; the
    ordinary column provides the across-sample baseline.
  \end{tablenotes}
  \end{threeparttable}
\end{table}
